%% kapa: exam=1 round=22 year=2011 period=2 form=A subject=accounting subject_ko=회계학 seq=33 number=33
%% source: Q-Net 감정평가사 기출문제 (cst00309) artlSeq=5036790 file=extracted/1008/5036790/2011년 제22회 감정평가사 1차시험 기출문제/2교시.hwp
%% conversion: kordoc 4.24.51 -> tools/gen1.py (auto); review=proofread-sonnet; flags=none
\begin{question}{33}
(주)서울은 단일제품을 생산ㆍ판매하고 있다. 20×1년 생산량은 500단위이고 판매량은 300단위이며, 원가자료는 다음과 같다.

{\small\centering
\begin{tabular}{|l|r|}
\hline
\multicolumn{1}{|c|}{항 목} & \multicolumn{1}{c|}{단위당 원가} \\
\hline
변동제조원가 & ￦23,000 \\
\hline
변동판매관리비 & 5,000 \\
\hline
\end{tabular}\par}

연간 고정제조간접원가는 ￦1,000,000이고 고정판매관리비는 ￦500,000이라면, 당기의 전부원가계산에 의한 영업이익과 변동원가계산에 의한 영업이익의 차이는? (단, 기초재고수량은 없으며, 단위당 판매가격은 ￦50,000임) \par
\begin{choices}
  \item ￦200,000
  \item ￦400,000
  \item ￦500,000
  \item ￦700,000
  \item ￦900,000
\end{choices}
\end{question}
\answer{2}
